\chapter{Совместный гессиан изотипической относительной кривизны}
\label{ch:v8-isotypic-relative-curvature-parent-hessian-gate}

\section{Проверка на прежнем физическом срезе}

Гладкий член предыдущего гейта имеет вид
\begin{equation}
 S_B(A)=\|AA^*B(A)-B(A)A^*A\|^2,
 \qquad
 B(A)=\Gamma_tA-A\Gamma_s.
 \label{eq:v8-isotypic-action-recall}
\end{equation}
Его гессиан в $A=0$ равен нулю, а в вакууме положительно
полуопределён. Поэтому на старом 27-мерном вещественном срезе действие
\begin{equation}
 S_{\rm VII}+\kappa S_B,
 \qquad \kappa\geq0,
 \label{eq:v8-old-slice-isotypic-extension}
\end{equation}
сохраняет исходную сигнатуру $(7,0,20)$ и вакуумную сигнатуру $(0,0,27)$.
Это проверено для восьми значений $\kappa$ от $0$ до $10^6$; минимальная
вакуумная мода не опускается ниже $5.6$.

Следовательно, новый член совместим с качественным переходом Тома VII. Но
эта проверка ещё не отвечает на вопрос о полном калибровочно-замкнутом модуле.

\section{Калибровочное замыкание меняет число запусковых направлений}

Полный модуль переноса имеет разложение
\begin{equation}
 \mathcal T_{15}=\mathcal I_5\oplus\mathcal H_{10},
 \label{eq:v8-incidence-heavy-complex-splitting}
\end{equation}
где $\mathcal I_5$ --- минимальная комплексная калибровочная орбита фона инцидентности,
а $\mathcal H_{10}$ --- тяжёлый модуль стрелок. На вещественной реализации
получается разложение $30=10+20$.

Ортогональный проектор $P_I$ на $\mathcal I_5$ коммутирует с полным
калибровочным представлением с остатком $1.74\cdot10^{-15}$. Поэтому минимальный
калибровочно-инвариантный запуск, содержащий прежний фон, уже не может состоять из
семи вещественных мод: он содержит пять комплексных, то есть десять
вещественных направлений.

Для комплексно-эрмитова квадратичного действия, сохраняющего скалярную фазу
$A\mapsto e^{i\theta}A$, вещественный индекс Морса чётен. Поэтому точное
число семь несовместимо с полным комплексный поле переноса без дополнительного
факторпространства вещественной структуры или нарушения этой фазовой симметрии.

\section{Что делает новая относительная кривизна на полном модуле}

Относительный Грам-гессиан на 30 вещественных направлениях имеет
\begin{equation}
 \operatorname{sig}H_{\rm Gram}^{0}=(30,0,0),
 \qquad
 \operatorname{sig}H_{\rm Gram}^{*}=(0,2,28).
 \label{eq:v8-full-transfer-gram-signatures-again}
\end{equation}
Гессиан $S_B$ имеет в вакууме ранг 12, но его ограничение на двумерное ядро
$H_{\rm Gram}^{*}$ равно нулю с точностью $3.2\cdot10^{-16}$. Значит,
изотипический член сам по себе не закрывает две оставшиеся плоские моды.

Рассмотрим минимальную калибровочно-инвариантную форму продолжения
\begin{align}
 H_0(m_I,m_H)
 &=-m_IP_I+m_H(1-P_I)+\frac12H_{\rm Gram}^{0},\label{eq:v8-full-origin-two-mass-completion}\\
 H_*(m_I,m_H)
 &=2m_IP_I+m_H(1-P_I)+\frac12H_{\rm Gram}^{*}+H_B.
 \label{eq:v8-full-vacuum-two-mass-completion}
\end{align}
Для представителя $(m_I,m_H)=(4,3.6)$ получено
\begin{equation}
 \operatorname{sig}H_0=(10,0,20),
 \qquad
 \operatorname{sig}H_*=(0,0,30),
 \qquad
 \lambda_{\min}(H_*)=5.53437\ldots.
 \label{eq:v8-full-transfer-qualitative-transition}
\end{equation}
Тем самым качественный калибровочно-замкнутый переход существует. Однако скан
показывает сигнатуры $(30,0,0)$, $(28,0,2)$ и $(10,0,20)$ при разных
положительных $m_I,m_H$. Эти два веса не выведены текущей архитектурой.

\begin{theorem}[Совместимость и минимальное калибровочное замыкание запуска]
Гладкая изотипическая относительная кривизна совместима со всем положительным
классом переходов Тома VII. На полном модуле переноса канонически выделяется
подмодуль инцидентности размерности пять над $\mathbb C$, поэтому минимальный
комплексно-самосопряжённый запуск содержит десять вещественных мод. Существует
открытая область двухмассовых завершений с переходом
$(10,0,20)\to(0,0,30)$.
\end{theorem}

\begin{nogocriterion}[Два массовых веса не являются выводом]
Нельзя объявлять формы \eqref{eq:v8-full-origin-two-mass-completion}--
\eqref{eq:v8-full-vacuum-two-mass-completion} единственным родительское действие.
Проекторы $P_I$ и $1-P_I$ выведены из калибровочного модуля, но их относительные
веса $m_I,m_H$ остаются свободными и меняют исходную сигнатуру. Нельзя также
требовать сохранения ровно семи отрицательных вещественных мод после перехода
к полному комплексному калибровочному замыканию без отдельного механизма вещественной структуры.
\end{nogocriterion}

\begin{protocol}
\begin{enumerate}
 \item старый срез: \textbf{$27$ вещественных измерений};
 \item совместимость нового члена с переходом $(7,0,20)\to(0,0,27)$:
 \textbf{пройдена};
 \item полный модуль переноса: \textbf{$15$ комплексных измерений, $30$ вещественных измерений};
 \item подмодуль инцидентности: \textbf{$5$ комплексных измерений, $10$ вещественных измерений};
 \item тяжёлый подмодуль: \textbf{$10$ комплексных измерений, $20$ вещественных измерений};
 \item проектор $P_I$: \textbf{калибровочно-инвариантен};
 \item $H_B$ поднимает Грам-ядро: \textbf{нет};
 \item представитель полного перехода: \textbf{$(10,0,20)\to(0,0,30)$};
 \item вакуумная щель представителя: \textbf{$5.53437$};
 \item единственные массовые веса: \textbf{не получены};
 \item полный родительское действие: \textbf{не получен};
 \item следующий гейт: \textbf{происхождение калибровочно-замкнутого рёберно-ходжева
 запуска}.
\end{enumerate}
\end{protocol}

Машинный сертификат сохранён в
\path{s2t_v8_isotypic_relative_curvature_parent_hessian_gate_results.json}.